% Anexo A: cómo reproducir los resultados. Anexo B: comprobaciones automáticas.

\section*{Anexo A. Reproducción de los resultados}

Todas las tablas y figuras del Capítulo~\ref{cap:resultados} se generan a partir de los lotes del
repositorio. Se requiere Python 3 con las dependencias de \texttt{requirements.txt}.

\subsection*{Preparar el entorno}

\begin{verbatim}
python3 -m venv venv
venv/bin/pip install -r requirements.txt
\end{verbatim}

\subsection*{Correr las pruebas de simulación}

Cada lote deja su salida en \texttt{reportes/}: el resumen estadístico, una fila por corrida y los
reportes de cada semilla. Entre paréntesis, la prueba de la Tabla~\ref{tab:experimentos} a la que
corresponde.

\begin{verbatim}
venv/bin/python main.py --batch lotes/tipos_paquete.json           # P3
venv/bin/python main.py --batch lotes/capitulo5_fallos.json        # P4
venv/bin/python main.py --batch lotes/capitulo5_fallos_movil.json  # P5
venv/bin/python main.py --batch lotes/movilidad.json               # P6
venv/bin/python main.py --batch lotes/movilidad_datos.json         # P7
venv/bin/python main.py --batch lotes/robustez.json                # P8
venv/bin/python main.py --batch lotes/generalidad.json             # P9
\end{verbatim}

\subsection*{Generar tablas y figuras}

La siguiente orden toma el lote más reciente de cada tipo y vuelve a escribir las tablas de
\texttt{Plantilla/MainMatter/Cap5/tablas/} y las figuras de \texttt{Plantilla/Images/Cap5/}. No
corre simulaciones.

\begin{verbatim}
venv/bin/python -m analysis.capitulo5
\end{verbatim}

\subsection*{Observar una prueba en la interfaz}

\begin{verbatim}
# Interfaz web local, con el escenario de la caída del puente (P1, P2 y P4)
venv/bin/python main.py --web --seed 1 \
    --config escenarios/fallos/fallo_puente.txt

# La misma corrida sin ventana, con el resumen en la terminal
venv/bin/python main.py --headless --duracion 240 --seed 1 \
    --config escenarios/fallos/fallo_puente.txt
\end{verbatim}

En la interfaz, la pestaña ``Tráfico'' permite elegir el tipo de paquete, la cantidad y el
intervalo; la pestaña ``Red'' muestra la matriz de conocimiento, y la tecla O muestra los OGM del
nodo seleccionado.

\subsection*{Compilar este documento}

\begin{verbatim}
cd Plantilla
pdflatex main && bibtex main && pdflatex main && pdflatex main
\end{verbatim}

El repositorio incluye una guía más amplia, con casos de uso y sus salidas esperadas, en
\texttt{docs/guia\_ejecucion.md}.

\section*{Anexo B. Comprobaciones automáticas del código}

Además de las pruebas de simulación, el repositorio incluye 390 comprobaciones automáticas escritas
con el módulo \texttt{unittest} de Python~\cite{Python}. Cada una ejecuta una parte pequeña del
programa y compara el resultado con el valor correcto; por ejemplo, que un OGM perdido baje el TQ,
que un enlace sin pérdidas lo mantenga en 1.0, o que una caída programada dé lo mismo que la misma
caída hecha a mano en la interfaz. Sirven para detectar si un cambio en el programa dañó algo que
ya funcionaba. La Tabla~\ref{tab:pruebas} muestra cómo se reparten. Se ejecutan con:

\begin{verbatim}
venv/bin/python -m unittest discover -s tests
\end{verbatim}

Las 390 pasan; en Linux se omite una, que necesita un componente de macOS.

\begin{table}[htbp!]
    \centering
    \caption{Comprobaciones automáticas del código, por área del simulador}
    \label{tab:pruebas}
    \footnotesize
    \begin{tabular}{|p{9.2cm}|p{2.6cm}|}
        \hline
        \textbf{Área} & \textbf{Comprobaciones} \\
        \hline
        Protocolo (TQ, elección de rutas) y medio de radiofrecuencia & 17 \\
        \hline
        Motor y nodos simulados & 57 \\
        \hline
        Lectura y revisión de escenarios & 26 \\
        \hline
        Fallos programables y adaptación de rutas & 30 \\
        \hline
        Envío de paquetes de datos y tipos de paquete & 38 \\
        \hline
        Movimiento descentralizado & 13 \\
        \hline
        Métricas de lo que cree cada nodo & 12 \\
        \hline
        Otros mapas, movimiento de usuarios e inundación & 14 \\
        \hline
        Reportes, lotes, uso por terminal y análisis & 82 \\
        \hline
        Interfaz web & 101 \\
        \hline
        \textbf{Total} & \textbf{390} \\
        \hline
    \end{tabular}
\end{table}
